\section{Conclusion}
Our experimental results demonstrate that NetTrack exhibits high universality, robustly adapting to diverse outdoor scenarios and various subtitle styles, including colored subtitles and those positioned centrally within video frames. In contrast, the Birds' Track approach, which relies on OpenCV, shows limited universality, achieving satisfactory performance only in simple scenes with specific subtitle characteristics, such as white subtitles located at the bottom of the video.

In terms of computational efficiency and resource utilization, the OpenCV-based method is highly efficient, operating exclusively on CPU with rapid execution times. Conversely, NetTrack requires significantly longer runtime and higher memory consumption, even when fine-tuned on CPU. Additionally, NetTrack demands substantial GPU resources and memory for training to obtain optimal weights, further highlighting its computational intensity.

However, this work has several limitations. Due to constraints in time and memory resources, our experiments were conducted solely on the VAL dataset, without leveraging the larger and more comprehensive train dataset. Moreover, we did not perform extensive ablation studies or provide comparative curves for ablation results, which restricts a deeper understanding of model components.

For future work, we plan to address these limitations by incorporating larger datasets and conducting thorough ablation experiments with detailed result visualizations. Furthermore, we believe that the Birds' Track architecture has potential for performance enhancement through optimization. We aim to explore more suitable OpenCV-based techniques to refine the Birds' Track framework, thereby improving its efficiency and applicability in broader contexts.